\chapter{Момент импульса, спин и статистика}

\section{От обычного вращения к собственному спину}

Формально речь пойдёт о моменте импульса, но на самом деле нас интересует спин. Спин — одно из самых простых и в то же время наименее интуитивных свойств элементарных частиц, и его структура вытекает из математики момента импульса.

Момент импульса — вектор. Его направление лежит вдоль оси вращения. Вдоль этой оси возможны два направления, поэтому нужна договорённость, позволяющая выбрать между ними: согните пальцы правой руки по направлению вращения, и большой палец укажет направление вектора момента импульса. Это определение, а не теорема.

Его величина зависит от того, какая масса присутствует, как эта масса распределена, каков размер системы и какова её угловая скорость. Распределение массы и размер учитываются в моменте инерции. Поэтому большой и маленький объекты могут нести один и тот же момент импульса, вращаясь с очень разными скоростями и обладая очень разными энергиями.

Обычный составной объект может обладать двумя концептуально различными видами момента импульса. Чашка, движущаяся по окружности вокруг внешней оси, имеет момент импульса, потому что её центр масс движется вокруг этой оси, даже если сама чашка не вращается ни вокруг какой оси, проходящей через неё. Это орбитальный момент импульса. Спин означает вращение вокруг внутренней оси.

Строгое определение использует систему отсчёта центра масс. Перейдите в систему отсчёта, в которой центр масс покоится, так что полный обычный импульс равен нулю. Всякий момент импульса, остающийся в этой системе, является спиновым моментом импульса. Объект при этом не обязан быть внутренне неподвижным; покоится только его центр масс.

Для баскетбольного мяча это различие не создаёт трудностей. Баскетбольный мяч в покое и тот же самый мяч, приведённый во вращение, очевидно распознаются как один и тот же объект в двух разных состояниях. Квантовая механика осложняет эту картину, потому что момент импульса дискретен. Нет непрерывной последовательности допустимых значений момента импульса, соединяющей состояние с нулевым моментом импульса с первым вращательным состоянием. Поэтому можно спросить, является ли первое вращательное состояние тем же самым объектом в движении или же вообще другим объектом.

Формулировка отчасти условна, но физический вопрос точен: сколько энергии требуется, чтобы достичь первого ненулевого вращательного состояния? Если зазор мал, возбуждённый объект всё ещё легко узнаваем. Если зазор огромен, возбуждение может радикально изменить объект. Попытка сообщить атому слишком сильное вращательное возбуждение может сорвать с него электроны и разрушить атом, а не оставить тот же самый атом спокойно вращающимся.

Теперь рассмотрим электрон, временно оставив в стороне его известный собственный спин. Можно ли заставить электрон вращаться вокруг внутренней оси? Математическая точка не имеет видимой внутренней структуры, которую можно было бы вращать, но эксперимент не устанавливает размер электрона, предоставляя нам маленькую жёсткую поверхность, вращение которой можно было бы наблюдать. Если электрон имеет чрезвычайно малый, но конечный размер, для перехода к его первому дополнительному вращательному возбуждению может просто понадобиться слишком много энергии. Поскольку энергия растёт при уменьшении размера, требуемая энергия может быть астрономической, возможно даже порядка энергии Планка. Эта оценка намеренно предварительна.\footnote{Короткий фрагмент, следующий непосредственно за этой оценкой, недостаточно вразумителен, чтобы на его основе можно было выполнить дальнейшую содержательную реконструкцию.}

Следовательно, величина спина элементарной частицы рассматривается как фиксированная характеристика. У электронов нет целого набора последовательно возрастающих величин собственного спина. Если бы возбуждение, во всех прочих отношениях похожее на электрон, имело другое полное значение спина, его классифицировали бы как другую частицу.

Фиксированная величина не фиксирует направление. Вообразим частицу под названием \emph{spintron}, определяемую, в частности, величиной её момента импульса. Если её спин направлен вдоль одной оси, инвариантность относительно вращений требует, чтобы частица того же типа могла иметь спин, направленный и вдоль другой оси. Ни одно пространственное направление не является внутренне предпочтительным. Величина спина остаётся фиксированной, тогда как его ориентация может меняться. Таким образом, в задаче есть две, по-видимому, противоречащие друг другу составляющие: квантованная величина и свободно меняющаяся ориентация. Алгебра момента импульса покажет, как они сосуществуют.

\begin{classroomqa}
\lecturetimestamp{00:12:25}
\audiencequestion{Почему для приведения маленькой частицы во вращение требуется энергия?}
\lecturerresponse{Для объекта массы \(M\) и характерного радиуса \(R\) момент инерции имеет порядок
\begin{equation}
I\sim cMR^2,
\end{equation}
где численный коэффициент \(c\) зависит от формы и от того, как распределена масса. Если вещество у внешнего края движется со скоростью \(v\), его обычный импульс имеет порядок \(Mv\), а момент импульса равен импульсу, умноженному на расстояние:
\begin{equation}
L\sim MvR.
\end{equation}
Эквивалентно, \(L=I\omega\) при \(v\sim\omega R\). Кинетическая энергия вращения равна
\begin{equation}
E_{\mathrm{rot}}=\frac{L^2}{2I}.
\end{equation}
При фиксированных \(M\) и \(L\),
\begin{equation}
E_{\mathrm{rot}}\sim\frac{L^2}{2cMR^2}\propto\frac{1}{R^2}.
\end{equation}
Чем меньше объект, тем больше энергии стоит данный момент импульса. Эта обратноквадратичная зависимость является классической. Дополнительный квантово-механический факт состоит в том, что момент импульса изменяется дискретными шагами, так что первый допустимый шаг не может быть сколь угодно малым.}
\end{classroomqa}

Математика, которой мы сейчас воспользуемся, абстрактна и поначалу довольно неинтуитивна. Её ценность в том, что из небольшого набора алгебраических правил следуют экспериментально проверяемые факты.

\section{Орбитальный момент импульса}

Начнём с одной частицы, движущейся по круговой орбите вокруг начала координат. Её момент импульса по модулю равен импульсу, умноженному на перпендикулярное расстояние до начала координат. В плоскости одно направление обхода можно назвать положительным, а противоположное — отрицательным. Знак фиксирует направление вектора момента импульса относительно выбранной оси.

Это ещё не спин. Частица имеет момент импульса потому, что у неё есть обычный импульс. Теперь поместим две частицы на противоположные концы диаметра и заставим их двигаться в одном и том же направлении вращения. Их обычные импульсы взаимно компенсируются, так что центр масс покоится, но их моменты импульса направлены одинаково и складываются. Система имеет нулевой полный импульс и ненулевой полный момент импульса: классическая модель спина.

Пусть положение частицы относительно начала координат задаётся выражением
\begin{equation}
\vec r=(x,y,z)=(x_1,x_2,x_3),
\end{equation}
а её импульс — выражением
\begin{equation}
\vec p=(p_x,p_y,p_z)=(p_1,p_2,p_3).
\end{equation}
Момент импульса определяется векторным произведением
\begin{equation}
\vec L=\vec r\times\vec p.
\end{equation}
При перестановке множителей вектор меняет знак: \(\vec p\times\vec r=-\vec r\times\vec p\). Стандартное соглашение — использовать \(\vec r\times\vec p\); его направление задаётся правилом правой руки.\footnote{\raggedright Кадр из видеолекции \nolinkurl{lecture_07_figure_02.png}, 00:18:49.\par}

\begin{figure}[h]
\centering
\includegraphics[width=0.72\textwidth]{lecture_07_figure_02.png}
\par
\begin{tikzpicture}[scale=1.0]
\draw[->] (-0.6,0) -- (4.8,0) node[right] {$x$};
\draw[->] (0,-0.6) -- (0,3.8) node[above] {$y$};
\coordinate (P) at (3.2,2.3);
\draw[->,thick] (0,0) -- (P) node[midway,above left] {$\vec r$};
\draw[dashed] (P) -- (3.2,0);
\draw[dashed] (P) -- (0,2.3);
\draw[->,thick] (P) -- ++(-1.0,1.25) node[above] {$\vec p$};
\node at (4.15,2.75) {$\vec L=\vec r\times\vec p$};
\end{tikzpicture}
\caption{Для частицы в положении \(\vec r\) с импульсом \(\vec p\) момент импульса \(\vec L=\vec r\times\vec p\) перпендикулярен плоскости, которую они образуют.}
\lectureframe{00:18:49}
\end{figure}

В декартовых компонентах
\begin{equation}
\begin{aligned}
L_x&=yp_z-zp_y,\qquad
L_y=zp_x-xp_z,\\
L_z&=xp_y-yp_x.
\end{aligned}
\end{equation}
Достаточно запомнить первую формулу и циклически переставлять обозначения
\[
x\longrightarrow y\longrightarrow z\longrightarrow x.
\]
Если циклически переставить и метку компоненты, и обозначения справа, следующая формула получится без изменения относительного знака. Эта мнемоника полезна именно потому, что знак легко случайно перепутать.

Для точечной частицы \(\vec r\times\vec p\) — это орбитальный момент импульса относительно выбранного начала координат. Для составной системы нужно сложить как векторы моменты импульса всех составляющих:
\begin{equation}
\vec L_{\mathrm{tot}}=\sum_a \vec r_a\times\vec p_a.
\end{equation}
Если вычислять эту сумму в системе отсчёта центра масс, она может остаться ненулевой, хотя
\[
\sum_a\vec p_a=0.
\]
Этот остающийся момент импульса и есть спин составной системы.

\section{Квантовые операторы и канонические коммутационные соотношения}

В квантовой механике наблюдаемые величины представляются операторами. Поэтому положение, импульс и каждая компонента момента импульса являются операторами. Их явная реализация в виде дифференциальных операторов здесь не требуется. Коммутационные соотношения содержат достаточно информации, чтобы вывести спектр момента импульса.

Для двух операторов \(A\) и \(B\)
\begin{equation}
[A,B]=AB-BA.
\end{equation}
Если две наблюдаемые коммутируют, они могут быть одновременно точно заданы. Если они не коммутируют, квантовая механика запрещает их одновременное точное задание.

Все три координаты положения можно измерять одновременно, поэтому компоненты положения коммутируют:
\begin{equation}
[x_i,x_j]=0.
\end{equation}
В частности, любой оператор коммутирует сам с собой. Было бы бессмысленно, если бы наблюдаемая величина была несовместима с другим измерением той же самой наблюдаемой.

Компоненты импульса также коммутируют:
\begin{equation}
[p_i,p_j]=0.
\end{equation}
Компонента координаты коммутирует с компонентой импульса в другом направлении, но не с соответствующей ей компонентой импульса. Эти три канонических утверждения можно записать вместе как
\begin{equation}
[x_i,x_j]=0,\qquad
[p_i,p_j]=0,\qquad
[x_i,p_j]=i\hbar\,\delta_{ij}.
\end{equation}
Так, например,
\[
[x,p_y]=0,\qquad [x,p_x]=i\hbar,
\]
а при перестановке порядка знак меняется:
\begin{equation}
[p_j,x_i]=-i\hbar\,\delta_{ij}.
\end{equation}

Момент импульса и постоянная Планка имеют одинаковую размерность. Момент импульса — это произведение расстояния на импульс, тогда как соотношение неопределённостей имеет в размерностном отношении вид \begin{equation}
\Delta x\,\Delta p\sim\hbar
\end{equation} с точностью до привычного численного множителя в точном неравенстве. Поэтому естественно, что момент импульса будет квантоваться в единицах \(\hbar\), хотя само квантование ещё не выведено.

\section{Вывод алгебры момента импульса}

Первая задача состоит в том, чтобы вычислить коммутатор двух компонент момента импульса. Из \[
L_x=yp_z-zp_y,\qquad
L_y=zp_x-xp_z,
\] получаем \begin{align}
[L_x,L_y]
&=[yp_z-zp_y,\;zp_x-xp_z] \\
&=[yp_z,zp_x]-[yp_z,xp_z] \notag\\
&\quad{}-[zp_y,zp_x]+[zp_y,xp_z].
\end{align}

\begin{classroomqa} \lecturetimestamp{00:31:41} \audiencequestion{ Можно ли запомнить формулы для компонент, циклически переставляя символы слева направо, а не циклически переходя по трём уравнениям сверху вниз? } \lecturerresponse{ Да. Допустима любая из этих циклических мнемоник. Единственная цель мнемоники — порождать компоненты в правильном циклическом порядке, отслеживая при этом знак. } \end{classroomqa}

Коммутатор можно понять, задавшись вопросом, можно ли множители в одном произведении переставить через множители в другом. В \[
[yp_z,xp_z],
\] оператор \(y\) коммутирует и с \(x\), и с \(p_z\), тогда как \(p_z\) коммутирует и с \(x\), и с самим собой. Никакого препятствия нет, поэтому этот член обращается в нуль. Те же рассуждения дают \[
[zp_y,zp_x]=0.
\] Остаются только две пары произведений.

Для первой из них единственное препятствие возникает, когда \(p_z\) переставляется через \(z\): \begin{align}
[yp_z,zp_x]
&=yp_zzp_x-zp_xyp_z \\
&=y(p_zz-zp_z)p_x \\
&=y[p_z,z]p_x \\
&=-i\hbar\,yp_x.
\end{align} Порядок в каноническом коммутаторе существен: \[
[p_z,z]=-[z,p_z]=-i\hbar.
\]

Для другой выжившей пары минусы в определениях \(L_x\) и \(L_y\) при перемножении дают плюс. Единственное препятствие возникает, когда \(z\) переставляется через \(p_z\): \begin{align}
[zp_y,xp_z]
&=zp_yxp_z-xp_zzp_y \\
&=x(zp_z-p_zz)p_y \\
&=x[z,p_z]p_y \\
&=i\hbar\,xp_y.
\end{align} Объединяя оба вклада, получаем \begin{align}
[L_x,L_y]
&=-i\hbar\,yp_x+i\hbar\,xp_y \\
&=i\hbar(xp_y-yp_x) \\
&=i\hbar L_z.
\end{align}

Это важное упрощение. Коммутатор двух компонент момента импульса даёт третью компоненту; никакие операторы координаты или импульса не остаются. Циклическая перестановка даёт \begin{equation}
\begin{aligned}
[L_x,L_y]&=i\hbar L_z,\qquad
[L_y,L_z]=i\hbar L_x,\\
[L_z,L_x]&=i\hbar L_y.
\end{aligned}
\end{equation} Эквивалентно, \begin{equation}
[L_i,L_j]=i\hbar\,\epsilon_{ijk}L_k.
\end{equation}

Алгебра замыкается на самих операторах момента импульса. После получения этих соотношений можно отложить в сторону \(x_i\) и \(p_i\), использовавшиеся для их вывода. Та же алгебра применима, когда полный момент импульса является суммой по многим составляющим. В частности, она применима в системе центра масс, где остающийся момент импульса — это спин. Поэтому коммутаторы момента импульса можно использовать абстрактно, не конструируя предварительно спин как движение маленьких составляющих частей.

Выбор осей \(x\), \(y\) и \(z\) произволен. Если перейти поворотом к любому другому набору взаимно перпендикулярных инерциальных декартовых осей, то штрихованные компоненты удовлетворяют тем же соотношениям: \begin{equation}
[L'_x,L'_y]=i\hbar L'_z,
\end{equation} вместе с его циклическими аналогами. Алгебра инвариантна относительно вращений, поэтому физические предсказания не могут зависеть от того, какие инерциальные оси были выбраны для вычисления.

\begin{classroomqa} \lecturetimestamp{00:39:09} \audiencequestion{ Не зависит ли спиновый момент импульса от переносов осей координат так же, как и от их поворотов? } \lecturerresponse{ Да. Внутренний спин не зависит от переноса. Орбитальный момент импульса относительно начала координат может измениться при переносе начала, но спин, остающийся в системе центра масс, от такого переноса не зависит. } \end{classroomqa}

Абстрактная алгебра может порождать бесконечно много тождеств, и большинство из них физически неинтересны. Её ценность проявляется, когда тождество становится измеримым фактом или что-то раскрывает об операционном смысле измерения. До этого места построение всё ещё главным образом алгебраическое. Следующий шаг выявляет дискретную структуру спина.

\begin{classroomqa} \lecturetimestamp{00:40:42} \audiencequestion{ Что если выбранные оси \(x\), \(y\) и \(z\) вращаются вместе с частицами? } \lecturerresponse{ Обсуждение ограничено инерциальными системами отсчёта. Вращающаяся система координат неинерциальна и вносит такие эффекты, как центробежная сила и сила Кориолиса. Момент импульса не инвариантен при переходе к совместно вращающейся системе отсчёта: система, которая в инерциальной системе явно вращается, в системе, вращающейся вместе с ней, может казаться покоящейся. Используемая здесь инвариантность относительно вращений относится к по-разному ориентированным инерциальным декартовым осям, а не к осям, которые вращаются во времени. } \end{classroomqa}

\section{Лестничные операторы и ось квантования}

Введём две комбинации \begin{equation}
L_+=L_x+iL_y,
\qquad
L_-=L_x-iL_y.
\end{equation} Они аналогичны операторам рождения и уничтожения гармонического осциллятора. Их роль состоит в повышении и понижении одной выбранной компоненты момента импульса.

В оси \(z\) нет ничего особенного, но если требуется указать компоненту, нужно выбрать какую-то ось. Выберем \(z\) в качестве оси квантования и запишем \begin{equation}
L_z\lvert m\rangle=m\hbar\,\lvert m\rangle.
\end{equation} Таким образом, \(m\) — это измеренная \(z\)-компонента в единицах \(\hbar\). Во время алгебраических вычислений часто удобно положить \(\hbar=1\); здесь она будет сохранена явно.

Коммутатор \([L_+,L_-]\) пропорционален \(L_z\), но для ближайшего рассуждения он не нужен. Сначала важно узнать коммутатор каждого лестничного оператора с \(L_z\). Из \[
[L_x,L_z]=-i\hbar L_y,
\qquad
[L_y,L_z]=i\hbar L_x,
\] получаем \begin{align}
[L_+,L_z]
&=[L_x+iL_y,L_z] \\
&=[L_x,L_z]+i[L_y,L_z] \\
&=-i\hbar L_y+i(i\hbar L_x) \\
&=-\hbar(L_x+iL_y) \\
&=-\hbar L_+.
\end{align} Если изменить знак у члена \(iL_y\), получим \begin{equation}
[L_-,L_z]=+\hbar L_-.
\end{equation} Эквивалентно, если поменять порядок в каждом коммутаторе на обратный, \begin{equation}
[L_z,L_+]=\hbar L_+,
\qquad
[L_z,L_-]=-\hbar L_-.
\end{equation}

В каждый момент времени точно определённой может быть только одна компонента момента импульса. Любые две различные компоненты не коммутируют, поэтому их нельзя задать одновременно. Выбранная точно определённая компонента задаёт ось квантования; выбор \(z\) — это вопрос удобства, а не физического предпочтения.

Теперь подействуем \(L_+\) на собственное состояние \(L_z\). Коммутатор можно переписать в виде \begin{equation}
L_zL_+=L_+L_z+\hbar L_+.
\end{equation} Следовательно, \begin{align}
L_z\bigl(L_+\lvert m\rangle\bigr)
&=\bigl(L_+L_z+\hbar L_+\bigr)\lvert m\rangle \\
&=m\hbar L_+\lvert m\rangle+\hbar L_+\lvert m\rangle \\
&=(m+1)\hbar\,L_+\lvert m\rangle.
\end{align} Если \(L_+\lvert m\rangle\) не равно нулю, то это собственное состояние \(L_z\), собственное значение которого увеличилось на одну единицу \(\hbar\). Точно так же \begin{equation}
L_z\bigl(L_-\lvert m\rangle\bigr)
=(m-1)\hbar\,L_-\lvert m\rangle.
\end{equation} Таким образом, \(L_-\) понижает \(m\) на единицу.

Это уже нетривиальный вывод. Каким бы ни оказался спектр \(L_z\), соседние значения отличаются на единицу: \begin{equation}
\ldots,\ m-2,\ m-1,\ m,\ m+1,\ m+2,\ \ldots
\end{equation} Концы спектра ещё не определены.

\section{Конечные мультиплеты и обозначение спина}

Для фиксированного типа частицы величина момента импульса фиксирована. Его момент импульса может переориентироваться, но неограниченное увеличение его величины в конце концов породило бы другой объект. Поэтому для этой частицы должна существовать максимально возможная \(z\)-компонента.

Как может существовать верхний конец у лестницы, если есть оператор повышения? Нужную модель даёт гармонический осциллятор. Его оператор уничтожения даёт нуль, когда доходит до нижнего состояния. Здесь оператор повышения должен давать нуль, когда доходит до верхнего состояния: \begin{equation}
L_+\lvert m_{\max}\rangle=0.
\end{equation} Утверждение, что \(L_+\) повышает собственное значение, верно только тогда, когда порождаемый им вектор ненулевой. Состояние наивысшего веса — как раз то исключительное состояние, на котором последовательность повышений обрывается.

Нижний конец фиксируется инвариантностью относительно вращений. Если момент импульса может быть максимально направлен вдоль \(+z\), то тот же физический объект можно повернуть так, чтобы он был максимально направлен вдоль \(-z\). Следовательно, \begin{equation}
m_{\min}=-m_{\max}.
\end{equation} Спектр симметричен относительно нуля и имеет единичный шаг. Возможны два семейства.

В первом семействе присутствует нуль: \begin{equation}
\begin{aligned}
m={}&-m_{\max},-m_{\max}+1,\ldots,-1,0,\\
&1,\ldots,m_{\max}-1,m_{\max}.
\end{aligned}
\end{equation} Тогда \(m_{\max}\) — целое число, а число состояний равно \begin{equation}
2m_{\max}+1.
\end{equation} Максимальное значение характеризует частицу и называется её спином. В общепринятом обозначении, \begin{equation}
s\equiv m_{\max},
\qquad
\lvert s,m\rangle\equiv\lvert m\rangle
\end{equation} внутри мультиплета с фиксированным спином.\footnote{\raggedright Кадр видеолекции \nolinkurl{lecture_07_figure_03.png}, 01:01:08.\par}

\begin{figure}[htbp] \centering
\includegraphics[width=0.70\textwidth]{lecture_07_figure_03.png} \caption{Условие наивысшего веса обрывает последовательность повышений, и допустимые значения \(m\) образуют конечную лестницу с единичным шагом.} \lectureframe{01:01:08} \end{figure}

Частица со спином нуль имеет только \(m=0\). Частица со спином один имеет \[
m=1,0,-1,
\] а частица со спином два имеет пять значений компоненты. Бозон Хиггса — пример частицы со спином нуль. Фотон и бозон \(Z\) имеют спин один. Гравитон, если он существует, имеет спин два. Составные объекты могут иметь спин три и гораздо больше: это возможно для ядер, а достаточно макроскопические вращающиеся объекты, такие как баскетбольные мячи и галактики, способны нести сколь угодно большой момент импульса. В обычном списке элементарных частиц отсутствует элементарная частица со спином три.

Другое симметричное семейство с единичным шагом не содержит нуля. Оно начинается с \[
m=\frac12,-\frac12
\] и продолжается с единичным шагом: \begin{equation}
\begin{aligned}
m={}&-m_{\max},-m_{\max}+1,\ldots,-\frac32,-\frac12,\\
&\frac12,\frac32,\ldots,m_{\max}.
\end{aligned}
\end{equation} Здесь \(m_{\max}\) — полуцелое число. Самый простой представитель имеет спин \(\tfrac12\), за ним следуют спин \(\tfrac32\), спин \(\tfrac52\) и так далее. Следовательно, мультиплеты момента импульса образуют семейства с целым и полуцелым спином.

\begin{classroomqa} \lecturetimestamp{01:03:46} \audiencequestion{ Есть ли максимум и у мультиплета с полуцелым спином? } \lecturerresponse{ Да. Лестница снова имеет наивысшее значение. Для частицы со спином \(\tfrac32\) максимум равен \(m_{\max}=\tfrac32\); для частицы со спином \(\tfrac52\) он равен \(m_{\max}=\tfrac52\). Разница лишь в том, что максимум и любое другое значение в мультиплете — полуцелые. } \end{classroomqa}

\begin{classroomqa}
\lecturetimestamp{01:04:04}
\audiencequestion{Разве не существует элементарных частиц со спином три и выше?}
\lecturerresponse{В рамках обсуждаемой релятивистской квантовой схемы в обычный список спинов элементарных частиц входят \(0,\tfrac12,1,\tfrac32,\) и \(2\). Составные объекты не ограничены этим диапазоном и могут иметь спин три или гораздо больший.}
\end{classroomqa}

\begin{classroomqa}
\lecturetimestamp{01:04:50}
\audiencequestion{Означают ли отрицательные значения на лестнице отрицательный спин?}
\lecturerresponse{Нет. Отрицательное \(m\) — это значение компоненты: оно означает, что момент импульса направлен противоположно выбранной оси. Спин частицы — это неотрицательное максимальное значение \(s=m_{\max}\).}
\end{classroomqa}

\begin{classroomqa}
\lecturetimestamp{01:05:10}
\audiencequestion{Если значения компоненты дискретны, то что происходит со всеми ориентациями между ними?}
\lecturerresponse{Существуют не просто две или три возможные пространственные ориентации. Промежуточные ориентации представляются квантовыми суперпозициями состояний базиса \(m\). Рассматриваемое ниже построение для спина-\(\tfrac12\) делает это явным.}
\end{classroomqa}

\begin{classroomqa}
\lecturetimestamp{01:05:42}
\audiencequestion{А как насчёт штопорообразного движения или волн с круговой поляризацией?}
\lecturerresponse{Электромагнитная волна с круговой поляризацией составлена из фотонов с круговой поляризацией, и эти фотоны несут момент импульса вдоль оси распространения. Слово ``штопор'' сочетает вращение с поступательным движением, поэтому для объекта, рассматриваемого в покое, оно не подходит. Для движущейся частицы существенной величиной является гелицитет: ориентация её спина относительно направления движения.}
\end{classroomqa}

\begin{classroomqa}
\lecturetimestamp{01:07:10}
\audiencequestion{Есть ли какая-нибудь известная элементарная частица со спином \(\tfrac32\)?}
\lecturerresponse{Ни одна элементарная частица со спином-\(\tfrac32\) не была непосредственно подтверждена. Суперсимметрия предполагает партнёра гравитона со спином-\(\tfrac32\), называемого гравитино, но его непосредственно не обнаружили.}
\end{classroomqa}

\begin{classroomqa}
\lecturetimestamp{01:08:01}
\audiencequestion{Перестали ли операторы момента импульса действовать на исходное пространство волновых функций или пространство кет-векторов?}
\lecturerresponse{Исходное спиновое пространство ещё не было задано. Эта алгебра используется, чтобы определить это пространство. Для фиксированного спина \(s\) его базисные векторы помечаются допустимыми значениями
\[
\lvert s,m\rangle,\qquad m=-s,-s+1,\ldots,s.
\]
Это внутреннее спиновое пространство отлично от обычного пространства пространственных волновых функций, на котором орбитальный момент импульса действует как \(\mathbf r\times\mathbf p\), хотя и то и другое реализует одну и ту же алгебру момента импульса.}
\end{classroomqa}

\section{Модуль момента импульса}

\begin{classroomqa}
\lecturetimestamp{01:09:32}
\audiencequestion{Наблюдаем ли модуль вектора момента импульса?}
\lecturerresponse{Да. Удобнее всего использовать его квадрат,
\[
L^2=L_x^2+L_y^2+L_z^2.
\]
Его значение в состоянии с наибольшим весом не равно просто \(s^2\hbar^2\). Даже когда \(L_z\) точно определён и максимален, некоммутирующие поперечные компоненты не могут обе одновременно быть точно равны нулю, поэтому они дают дополнительный квантовый вклад.}
\end{classroomqa}

Чтобы вычислить этот вклад, перемножим лестничные операторы в таком порядке, при котором повышающий оператор сначала действует на кет:
\begin{align}
L_-L_+
&=(L_x-iL_y)(L_x+iL_y) \\
&=L_x^2+L_y^2+iL_xL_y-iL_yL_x \\
&=L_x^2+L_y^2+i[L_x,L_y] \\
&=L_x^2+L_y^2-\hbar L_z.
\end{align}
На последнем шаге использовано \([L_x,L_y]=i\hbar L_z\). Поэтому
\begin{equation}
L_x^2+L_y^2=L_-L_++\hbar L_z,
\end{equation}
и, следовательно,
\begin{equation}
L^2=L_-L_++L_z^2+\hbar L_z.
\end{equation}

Применим это тождество к состоянию с наибольшим весом \(\lvert s,s\rangle=\lvert m_{\max}\rangle\):
\begin{align}
L^2\lvert s,s\rangle
&=\left(L_-L_++L_z^2+\hbar L_z\right)\lvert s,s\rangle \\
&=L_-L_+\lvert s,s\rangle
+s^2\hbar^2\lvert s,s\rangle
+s\hbar^2\lvert s,s\rangle.
\end{align}
Первый член обращается в нуль, потому что \(L_+\lvert s,s\rangle=0\). Поэтому
\begin{equation}
L^2\lvert s,s\rangle
=s(s+1)\hbar^2\lvert s,s\rangle.
\end{equation}
Дополнительное \(+s\) — это квантовая поправка, связанная с неизбежными поперечными флуктуациями. Классически ожидалось бы только \(s^2\hbar^2\).

Первые несколько значений делают это различие наглядным:
\begin{align}
s=0 &: \qquad L^2=0,\\
s=1 &: \qquad L^2=2\hbar^2,\\
s=\frac12 &: \qquad L^2=\frac34\hbar^2.
\end{align}
Для \(s=1000\) безразмерный множитель равен \(1000\times1001\), что мало отличается от \(1000^2\) при столь большом значении момента импульса. Вообще,
\begin{equation}
s(s+1)=s^2\left(1+\frac1s\right),
\end{equation}
поэтому при очень больших \(s\) относительная квантовая поправка становится пренебрежимо малой. Так и возникает классический результат \(s^2\).

Каждое \(m\)-состояние в мультиплете с фиксированным спином имеет одно и то же собственное значение \(L^2\). Изменение \(m\) меняет компоненту вдоль выбранной оси, а не полный модуль. Это квантовый аналог поворота вектора момента импульса без изменения его длины.

\begin{classroomqa}
\lecturetimestamp{01:16:43}
\audiencequestion{Можно ли одновременно знать \(L^2\) и \(z\)-компоненту?}
\lecturerresponse{Да. Полный квадрат коммутирует с любой компонентой момента импульса:
\[
[L^2,L_i]=0,\qquad i=x,y,z.
\]
Следовательно, \(L^2\) и одну выбранную компоненту, такую как \(L_z\), можно задавать одновременно, хотя две разные компоненты нельзя. У любого электрона, например, \(L^2=\tfrac34\hbar^2\), независимо от его ориентации.}
\end{classroomqa}

\begin{classroomqa}
\lecturetimestamp{01:17:41}
\audiencequestion{Равно ли \(L^2\) нулю при спине 1, или только при нулевом спине?}
\lecturerresponse{Только у нулевого спина \(L^2=0\). Для спина 1 получаем
\[
L^2=1(1+1)\hbar^2=2\hbar^2,
\]
а для спина 3
\[
L^2=3(3+1)\hbar^2=12\hbar^2.
\]}
\end{classroomqa}

\begin{classroomqa}
\lecturetimestamp{01:18:13}
\audiencequestion{Почему каждый уровень \(m\) в мультиплете имеет одно и то же собственное значение \(L^2\)?}
\lecturerresponse{Начнём с состояния с наибольшим весом и будем понижать его. Поскольку \(L^2\) коммутирует с каждой компонентой, оно коммутирует и с \(L_-\):
\[
[L^2,L_-]=0.
\]
Если
\[
L^2\lvert s,s\rangle=s(s+1)\hbar^2\lvert s,s\rangle,
\]
то
\begin{align}
L^2\bigl(L_-\lvert s,s\rangle\bigr)
&=L_-L^2\lvert s,s\rangle \\
&=s(s+1)\hbar^2\,L_-\lvert s,s\rangle.
\end{align}
Следовательно, следующее состояние ниже имеет то же собственное значение \(L^2\). Повторение этого рассуждения сохраняет это собственное значение по всей лестнице до самого низа.}
\end{classroomqa}

Следовательно, дискретность значений \(m\) не означает, что существует только дискретный набор пространственных ориентаций. Инвариантность относительно вращений по-прежнему требует состояний, соответствующих повёрнутым направлениям. Для спина \(\tfrac12\) эти направления будут возникать как суперпозиции двух базисных состояний \(L_z\).

\section{Состояния со спином-\texorpdfstring{\(\tfrac12\)}{1/2} и ориентация}

Инвариантность относительно вращений требует большего, чем одни лишь дискретные собственные состояния \(L_z\). Спин может быть направлен вдоль любой оси, хотя измерение его компоненты вдоль выбранной оси даёт только допустимые дискретные значения. Направления между собственными состояниями для оси \(z\) представляются суперпозициями.

Для частицы со спином-\(\tfrac12\) ограничимся спиновыми степенями свободы. У электрона есть и орбитальные состояния, так что всего у него, конечно, не только два состояния; двумерно лишь его спиновое пространство. Выберем базис по оси \(z\)
\begin{equation}
L_z\lvert+\rangle=\frac{\hbar}{2}\lvert+\rangle,
\qquad
L_z\lvert-\rangle=-\frac{\hbar}{2}\lvert-\rangle.
\end{equation}
Общее спиновое состояние имеет вид
\begin{equation}
\lvert\psi\rangle
=
\alpha\lvert+\rangle+\beta\lvert-\rangle,
\end{equation}
где \(\alpha\) и \(\beta\) — комплексные числа. Измерение \(L_z\) даёт \(+\hbar/2\) с вероятностью \(|\alpha|^2\) и \(-\hbar/2\) с вероятностью \(|\beta|^2\), поэтому условие нормировки имеет вид
\begin{equation}
|\alpha|^2+|\beta|^2=1.
\end{equation}

Общая фаза не оказывает физического эффекта:
\begin{equation}
(\alpha,\beta)\longrightarrow
e^{i\chi}(\alpha,\beta).
\end{equation}
Две комплексные амплитуды изначально содержат четыре вещественных параметра. Нормировка устраняет один, а физически несущественная общая фаза устраняет ещё один, так что остаются два физических вещественных параметра. Это ровно то число, которое нужно, чтобы задать направление полярным и азимутальным углами. Изменение \(\alpha\) и \(\beta\) поворачивает ориентацию спина.

Измеренное значение компоненты и её среднее значение — разные вещи. Среднее значение — это среднее по ансамблю. Подготовьте много электронов в точно одном и том же состоянии, измерьте \(L_z\) по одному разу для каждого электрона, сложите результаты и разделите на число экспериментов. Каждый отдельный результат по-прежнему равен либо \(+\hbar/2\), либо \(-\hbar/2\), но среднее может принимать любое значение между ними:
\begin{equation}
\langle L_z\rangle
=
\frac{\hbar}{2}\bigl(|\alpha|^2-|\beta|^2\bigr).
\end{equation}

Равные амплитуды состояний вверх и вниз дают спин, направленный вдоль одной из ориентаций оси \(x\). Нормированные представители таковы\footnote{Отношения коэффициентов определяют направления. Множители \(1/\sqrt2\) обеспечивают нормировку, требуемую условием \( |\alpha|^2+|\beta|^2=1\). Относительные фазы \(+i\) и \(-i\) задают противоположные ориентации оси \(y\); то, какой из знаков соответствует положительному \(y\), зависит от соглашения о направлении оси.}
\begin{equation}
\lvert x;+\rangle
=
\frac{\lvert+\rangle+\lvert-\rangle}{\sqrt2},
\qquad
\lvert x;-\rangle
=
\frac{\lvert+\rangle-\lvert-\rangle}{\sqrt2}.
\end{equation}
Для любого из этих состояний измерение вдоль \(z\) даёт два возможных результата с равной вероятностью, и поэтому
\begin{equation}
\langle L_z\rangle
=
\frac12\left(\frac{\hbar}{2}\right)
+
\frac12\left(-\frac{\hbar}{2}\right)
=0.
\end{equation}
Относительные знаки плюс и минус задают противоположные направления вдоль оси \(x\). Аналогично,
\begin{equation}
\frac{\lvert+\rangle+i\lvert-\rangle}{\sqrt2},
\qquad
\frac{\lvert+\rangle-i\lvert-\rangle}{\sqrt2}
\end{equation}
задают две противоположные ориентации оси \(y\). Произвольные нормированные значения \(\alpha\) и \(\beta\) описывают произвольные оси. Сказать, что спин направлен вдоль некоторой оси, значит, что измерение компоненты вдоль этой оси даёт определённый результат. Относительно другой оси то же состояние, вообще говоря, является суперпозицией.

Индивидуальные результаты измерения остаются квантованными. Непрерывно меняются именно их средние по ансамблю, и они воспроизводят привычную геометрию классического вектора.

\section{Спин, статистика и составные частицы}

Масса, электрический заряд и спин относятся к основным свойствам, используемым для характеристики частицы. Спин также связан со статистикой. Соотношение спин-статистика — теорема релятивистской квантовой механики, хотя здесь оно доказываться не будет:
\begin{equation}
\begin{aligned}
\text{half-integer spin} &\quad\Longleftrightarrow\quad \text{fermion},\\
\text{integer spin} &\quad\Longleftrightarrow\quad \text{boson}.
\end{aligned}
\end{equation}
Два тождественных фермиона не могут занимать одно и то же полное квантовое состояние. Тождественные бозоны могут занимать одно и то же состояние, и большие их количества в одном состоянии могут коллективно вести себя как классическая волна. Поэтому частицы со спином \(\tfrac12,\tfrac32,\tfrac52,\ldots\) являются фермионами, а частицы со спином \(0,1,2,\ldots\) — бозонами.

Это правило сразу распространяется и на составные системы. Сложение двух полуцелых моментов импульса дает целый момент импульса. В более общем случае четное число фермионных составляющих образует бозонную составную систему, тогда как нечетное — фермионную составную систему. Добавление бозонных составляющих не меняет этого правила четности.

Обычный атом водорода содержит протон со спином \(\tfrac12\) и электрон со спином \(\tfrac12\). Поэтому в нем четное число фермионов, и в целом он является бозоном. Дейтрон содержит протон и нейтрон, оба фермионы, и тоже является бозоном. В атоме дейтерия к дейтрону добавляется электрон, так что на этом уровне описания получается три фермионные составляющие, следовательно, дейтерий — фермион.

Частицы и античастицы имеют одинаковые массу и спин и противоположные электрические заряды. Позитроний — это электрон и позитрон, связанные друг с другом. Он содержит два фермиона и поэтому является бозоном.

\begin{classroomqa}
\lecturetimestamp{01:33:08}
\audiencequestion{Проходят ли бозоны сквозь вещество, и связана ли сверхпроводимость с тем, что спаренные электроны ведут себя как бозоны?}
\lecturerresponse{Существенное свойство состоит не просто в способности проходить сквозь вещество; оно состоит в способности бозонов занимать одно и то же состояние. Более простой пример дает сверхтекучий гелий. Обычный атом гелия-\(4\) содержит два протона, два нейтрона и два электрона, то есть четное число фермионных составляющих, поэтому такой атом является бозоном. Многие атомы гелия могут занимать одно квантовое состояние и двигаться коллективно. Тогда сверхтекучую жидкость можно приблизительно представить как классическую волну, образующуюся, когда много атомов помещены в одно и то же состояние. Эта картина неточна, но передает нужную здесь мысль. Электрон в металле — фермион и не может находиться в одном и том же полном состоянии с другим электроном. Пары электронов, однако, являются бозонными составными системами; в этом и состоит связь со сверхпроводимостью. Здесь возникает парадокс: если составляющие электроны не могут занимать одно и то же состояние, то как несколько составных пар могут занимать одно состояние? Разрешение этого парадокса требует более аккуратного рассмотрения составных состояний и откладывается.}
\end{classroomqa}

\begin{classroomqa}
\lecturetimestamp{01:35:46}
\audiencequestion{Поскольку протоны и нейтроны составные, почему у них полуцелый спин?}
\lecturerresponse{Каждый из них состоит из нечетного числа кварков со спином \(\tfrac12\). Три фермионные составляющие образуют фермионную составную систему, поэтому протоны и нейтроны могут иметь полуцелый спин.}
\end{classroomqa}

\section{Принцип запрета Паули и атомные орбитали}

Непосредственную мотивацию для принципа запрета Паули дает периодическая таблица. Правило не состоит в том, что два электрона вообще никогда не могут иметь одно и то же орбитальное движение. Оно состоит в том, что два электрона не могут занимать одно и то же полное квантовое состояние.

Это различие важно даже в самых простых атомах. Если, начиная с водорода, увеличить заряд ядра и добавить второй электрон, то два электрона могут занимать одну и ту же пространственную орбиталь, но их спиновые состояния должны различаться. Один может находиться в состоянии со спином вверх, а другой — со спином вниз относительно выбранной оси. Когда два электрона занимают одну орбиталь, их спины образуют спиновое синглетное состояние, так что их суммарный спин взаимно компенсируется. Одну пространственную орбиталь могут занимать не более двух электронов, потому что имеется только два независимых спиновых состояния.

\begin{classroomqa}
\lecturetimestamp{01:37:53}
\audiencequestion{Можно ли представлять противоположные спины двух электронов на одной орбитали как сцепленные шестерни?}
\lecturerresponse{Как грубая наглядная картина — да: сцепленные шестерни вращаются в противоположных направлениях. Это лишь аналогия для противоположных направлений спина, а не буквальная механическая модель электронов.}
\end{classroomqa}

Слово ``орбиталь'' может употребляться на двух разных уровнях. Состояние орбитального движения в широком смысле — это состояние, получаемое при решении уравнения Шрёдингера без учета спина. Оно содержит больше информации, чем просто то, какую оболочку занимает электрон.

\begin{classroomqa}
\lecturetimestamp{01:38:09}
\audiencequestion{Означает ли орбиталь нечто большее, чем просто то, какую оболочку занимает электрон?}
\lecturerresponse{Да. Здесь под орбитальным движением понимается шрёдингеровское состояние электрона без учета спина. Оно включает главное квантовое число \(n\), которое несет радиальную информацию и приблизительно задает расстояние электрона от ядра, а также квантовые числа орбитального момента импульса. В спектроскопической нотации, используемой химиками, \(S\), \(P\), \(D\) и \(F\) соответствуют орбитальному моменту импульса \(l=0,1,2,3\) соответственно.}
\end{classroomqa}

Таким образом, главное квантовое число и орбитальный момент импульса описывают разные стороны движения. Обозначения \(S,P,D,F\) относятся к значению \(l\), тогда как \(n\) дает отдельную радиальную информацию. Правило Паули состоит в том, что у двух электронов не могут совпадать все их квантовые числа, включая спиновое квантовое число.

\begin{classroomqa}
\lecturetimestamp{01:40:42}
\audiencequestion{Позволяют ли квантовые числа определить, сколько электронов может занимать каждую оболочку в атоме?}
\lecturerresponse{Да. Они определяют максимально возможное число, хотя не то, сколько электронов у данного атома действительно находится в этой оболочке. Общая задача подсчета по оболочкам здесь рассматриваться не будет.}
\end{classroomqa}

Для фиксированного орбитального момента импульса \(l\) квантовое число его компоненты имеет
\begin{equation}
2l+1
\end{equation}
возможных значений. Это тот же самый подсчет для момента импульса, который уже использовался для спина. Состояние \(S\) имеет \(l=0\), поэтому ему соответствует одно орбитальное состояние. Состояние \(P\) имеет \(l=1\), поэтому ему соответствуют три орбитальных состояния. Состояние \(D\) имеет \(l=2\), поэтому ему соответствуют пять.

Спин удваивает число состояний для каждой орбитали, потому что каждая пространственная орбиталь может содержать по одному электрону в каждом из двух противоположных спиновых состояний.

\begin{classroomqa}
\lecturetimestamp{01:41:54}
\audiencequestion{Вмещает ли состояние \(S\) два электрона, а три состояния \(P\) — еще шесть?}
\lecturerresponse{Да. Единственное орбитальное состояние \(S\) вмещает не более двух электронов. Три орбитальных состояния \(P\) вмещают по два электрона каждое, что дает еще шесть. Следовательно, рассматриваемые состояния \(S\) и \(P\) в сумме вмещают \(2+6=8\) электронов.}
\end{classroomqa}

Та же структура момента импульса, которая упорядочивает атомные орбитали, вновь появляется, с некоторыми отличиями, когда кварки движутся в связанных системах и обращаются друг вокруг друга.